% Môn: Vật lý
% Lớp: 12
% Chương: Chương I. Vật lí nhiệt
% Bài: L12C1 Bài 1. Cấu trúc của chất. Sự chuyển thể
% Số câu: 194
% App đọc file này trực tiếp — sửa rồi Commit trên GitHub, bấm Đồng bộ GitHub .tex

% L12C1 Bài 1. Cấu trúc của chất. Sự chuyển thể
% ID: L12C1B1-TN-42
% Mức: NB
% ID: L12C1B1-91-TN


% ID: L12C1B1-TN-4
% Mức: NB
% ID: L12C1B1-97-TN
% ID: L12C1B1-97-TN
% Mức: NB
% ID: L12C1B1-97-TN
% Mức: NB
% ID: L12C1B1-97-TN


% ID: L12C1B1-TN-3
% Mức: NB
% ID: L12C1B1-100-TN
% ID: L12C1B1-100-TN
% Mức: NB
% ID: L12C1B1-100-TN
% Mức: NB
% ID: L12C1B1-100-TN
% ID: L12C1B1-100-TN
% Mức: NB
% ID: L12C1B1-100-TN
% Mức: NB
% ID: L12C1B1-100-TN
% ID: L12C1B1-100-TN
% Mức: NB
\dangbt{Mô hình động học phân tử về cấu tạo chất và chuyển động Brown}
\begin{ex}
Điều nào sau đây là đúng khi nói về mô hình động học phân tử?
\choice
{Các hạt phân tử cấu tạo nên các chất luôn đứng yên tại một vị trí cố định.}
{\True Các hạt phân tử cấu tạo nên các chất luôn chuyển động không ngừng.}
{Các hạt phân tử cấu tạo nên các chất có thể chuyển động hoặc đứng yên tùy vào đó là chất rắn, chất lỏng hay chất khí.}
{Các hạt phân tử cấu tạo nên chất lỏng và chất khí thì chuyển động, các hạt cấu tạo nên chất rắn thì đứng yên.}
\loigiai{Theo mô hình động học phân tử, các hạt phân tử cấu tạo nên các chất luôn chuyển động không ngừng.}
\end{ex}

\dangbt{Mô hình động học phân tử về cấu tạo chất và chuyển động Brown}
\begin{ex}
Theo mô hình động học phân tử nhiệt độ của vật cao hay thấp là do
\choice
{\True sự chuyển động nhanh hay chậm của các phân tử cấu tạo nên vật.}
{số phân tử cấu tạo nên vật là nhiều hay ít.}
{mật độ các phân tử trên một đơn vị thể tích là lớn hay nhỏ.}
{khối lượng của các phân tử là nặng hay nhẹ.}
\loigiai{Theo mô hình động học phân tử, nhiệt độ của vật cao hay thấp là do sự chuyển động nhanh hay chậm của các phân tử cấu tạo nên vật.}
\end{ex}

\dangbt{Phân biệt đặc điểm cấu trúc phân tử của chất ở thể rắn, lỏng, khí}
\begin{ex}
Với mô hình động học phân tử, sự khác biệt về cấu trúc của chất rắn, chất lỏng, chất khí là do sự khác biệt về
\choice
{thành phần các phân tử cấu tạo của mỗi chất.}
{\True độ lớn của lực tương tác giữa các phân tử trong mỗi chất.}
{số lượng phân tử cấu tạo nên mỗi chất.}
{kích thước của các phân tử cấu tạo của mỗi chất.}
\loigiai{Sự khác biệt về cấu trúc của chất rắn, chất lỏng, chất khí là do sự khác biệt về độ lớn của lực tương tác giữa các phân tử trong mỗi chất.}
\end{ex}

\dangbt{Mô hình động học phân tử về cấu tạo chất và chuyển động Brown}
\begin{ex}
Mô hình động học phân tử được xây dựng dựa trên quan điểm nào về cấu tạo chất?
\choice
{Dựa trên quan điểm chất có cấu tạo liên tục.}
{\True Dựa trên quan điểm chất có cấu tạo gián đoạn.}
{Dựa trên cả hai quan điểm là chất có cấu tạo liên tục và gián đoạn.}
{Dựa trên quan điểm về sự bền vững của các phân tử.}
\loigiai{Mô hình động học phân tử được xây dựng dựa trên quan điểm chất có cấu tạo gián đoạn.}
\end{ex}

\dangbt{Mô hình động học phân tử về cấu tạo chất và chuyển động Brown}
\begin{ex}
Câu nào sau đây nói về chuyển động của phân tử là không đúng?
\choice
{\True Chuyển động của phân tử là do lực tương tác phân tử gây ra.}
{Các phân tử chuyển động không ngừng.}
{Các phân tử chuyển động chậm thì nhiệt độ của vật thấp.}
{Các phân tử khí chuyển động theo đường thẳng giữa hai lần va chạm.} 
\loigiai{Chuyển động của phân tử là chuyển động nhiệt, không phải do lực tương tác phân tử gây ra.}
\end{ex}

\dangbt{Mô hình động học phân tử về cấu tạo chất và chuyển động Brown}
\begin{ex}
Câu nào sau đây nói về lực tương tác phân tử là không đúng?
\choice
{Lực tương tác phân tử chỉ đáng kể khi các phân tử ở rất gần nhau.}
{Lực hút phân tử có thể lớn hơn lực đẩy phân tử.}
{\True Lực hút phân tử không thể lớn hơn lực đẩy phân tử.}
{Lực hút phân tử có thể bằng lực đẩy phân tử.}
\loigiai{Lực hút phân tử có thể lớn hơn lực đẩy phân tử (khi khoảng cách lớn hơn một giá trị nhất định) hoặc nhỏ hơn lực đẩy phân tử (khi khoảng cách rất nhỏ).}
\end{ex}

\dangbt{Phân biệt đặc điểm cấu trúc phân tử của chất ở thể rắn, lỏng, khí}
\begin{ex}
Khi nói về khoảng cách trung bình giữa các phân tử trong chất rắn, chất lỏng, chất khí. Kết luận nào sau đây là đúng?
\choice
{Khoảng cách giữa các phân tử trong chất lỏng xa hơn so với các phân tử trong chất khí.}
{Khoảng cách giữa các phân tử trong chất rắn xa hơn so với các phân tử trong chất lỏng.}
{\True Khoảng cách giữa các phân tử trong chất lỏng gần hơn so với các phân tử trong chất khí.}
{Khoảng cách giữa các phân tử trong chất lỏng xa hơn so với các phân tử trong chất khí.}
\loigiai{Khoảng cách giữa các phân tử trong chất rắn là nhỏ nhất, sau đó đến chất lỏng và lớn nhất là chất khí. Vậy khoảng cách giữa các phân tử trong chất lỏng gần hơn so với các phân tử trong chất khí.}
\end{ex}

\dangbt{Mô hình động học phân tử về cấu tạo chất và chuyển động Brown}
\begin{ex}
Khi khoảng cách giữa các phân tử rất nhỏ, thì giữa các phân tử
\choice
{chỉ có lực hút.}
{chỉ có lực đẩy.}
{\True có cả lực hút và lực đẩy, nhưng lực đẩy lớn hơn lực hút.}
{có cả lực hút và lực đẩy, nhưng lực đẩy nhỏ hơn lực hút.}
\loigiai{Khi khoảng cách giữa các phân tử rất nhỏ, thì giữa các phân tử có cả lực hút và lực đẩy, nhưng lực đẩy lớn hơn lực hút.}
\end{ex}

\dangbt{Phân biệt đặc điểm cấu trúc phân tử của chất ở thể rắn, lỏng, khí}
\begin{ex}
Tính chất không phải là của phân tử của vật chất ở thể khí là
\choice
{chuyển động hỗn loạn.}
{chuyển động không ngừng.}
{chuyển động hỗn loạn và không ngừng.}
{\True chuyển động hỗn loạn xung quanh các vị trí cân bằng cố định.}
\loigiai{Tính chất không phải là của phân tử của vật chất ở thể khí là chuyển động hỗn loạn xung quanh các vị trí cân bằng cố định, đây là tính chất của phân tử ở thể rắn.}
\end{ex}

\dangbt{Phân biệt đặc điểm cấu trúc phân tử của chất ở thể rắn, lỏng, khí}
\begin{ex}
Tính chất nào sau đây không phải là tính chất của chất ở thể khí?
\choice
{\True Có hình dạng và thể tích riêng.}
{Có các phân tử chuyển động hoàn toàn hỗn độn.}
{Có thể nén được dễ dàng.}
{Có lực tương tác phân tử nhỏ hơn lực tương tác phân tử ở thể rắn và thể lỏng.}
\loigiai{Chất khí không có hình dạng và thể tích riêng, chúng chiếm toàn bộ thể tích bình chứa.}
\end{ex}

\dangbt{Phân biệt đặc điểm cấu trúc phân tử của chất ở thể rắn, lỏng, khí}
\begin{ex}
Phát biểu nào sau đây là sai khi nói về chất khí?
\choice
{Lực tương tác giữa các nguyên tử, phân tử rất yếu.}
{\True Các phân tử khí ở rất gần nhau.}
{Chất khí không có hình dạng và thể tích riêng.}
{Chất khí luôn chiếm toàn bộ thể tích bình chứa và có thể nén được dễ dàng.}
\loigiai{Các phân tử khí ở rất xa nhau, khoảng cách giữa chúng lớn hơn nhiều so với kích thước của phân tử.}
\end{ex}

\dangbt{Mô hình động học phân tử về cấu tạo chất và chuyển động Brown}
\begin{ex}
Phát biểu nào sau đây nói về chuyển động của phân tử là không đúng?
\choice
{\True Chuyển động của phân tử là do lực tương tác phân tử gây ra.}
{Các phân tử chuyển động không ngừng.}
{Các phân tử chuyển động càng nhanh thì nhiệt độ của vật càng cao.}
{Các phân tử khí lí tưởng chuyển động theo đường thẳng giữa hai vật va chạm.}
\loigiai{Chuyển động của phân tử là chuyển động nhiệt, không phải do lực tương tác phân tử gây ra.}
\end{ex}

\dangbt{Phân biệt đặc điểm cấu trúc phân tử của chất ở thể rắn, lỏng, khí}
\begin{ex}
Các nguyên tử, phân tử trong chất rắn
\choice
{\True nằm ở những vị trí xác định và chỉ có thể dao động xung quanh các vị trí cân bằng này.}
{nằm ở những vị trí cố định.}
{không có vị trí cố định mà luôn thay đổi.}
{nằm ở những vị trí cố định, sau một thời gian nào đó chúng lại chuyển sang một vị trí cố định khác.}
\loigiai{Trong chất rắn, các nguyên tử, phân tử nằm ở những vị trí xác định và chỉ có thể dao động xung quanh các vị trí cân bằng này.}
\end{ex}

\dangbt{Mô hình động học phân tử về cấu tạo chất và chuyển động Brown}
\begin{ex}
Các phân tử khí ở áp suất thấp và nhiệt độ tiêu chuẩn có các tính chất là
\choice
{chuyển động không ngừng và coi như chất điểm.}
{coi như chất điểm và tương tác hút hoặc đẩy với nhau.}
{chuyển động không ngừng và tương tác hút hoặc đẩy với nhau.}
{\True Chuyển động không ngừng, coi như chất điểm, và tương tác hút hoặc đẩy với nhau.}
\loigiai{Vì ở áp suất thấp, ta có thể coi khí thực gần đúng là khí lý tưởng nên chúng chỉ tương tác với nhau khi va chạm.}
\end{ex}

\dangbt{Mô hình động học phân tử về cấu tạo chất và chuyển động Brown}
\begin{ex}
Tính chất nào sau đây không phải là của phân tử?
\choice
{Chuyển động không ngừng.}
{Giữa các phân tử có khoảng cách.}
{\True Có lúc đứng yên, có lúc chuyển động.}
{Chuyển động càng nhanh thì nhiệt độ của vật càng cao.}
\loigiai{Các phân tử luôn chuyển động không ngừng, không có lúc đứng yên.}
\end{ex}

\dangbt{Phân biệt đặc điểm cấu trúc phân tử của chất ở thể rắn, lỏng, khí}
\begin{ex}
Trong điều kiện chuẩn về nhiệt độ và áp suất thì
\choice
{\True số phân tử trong một đơn vị thể tích của các chất khí khác nhau là như nhau.}
{các phân tử của các chất khí khác nhau chuyển động với vận tốc như nhau.}
{khoảng cách giữa các phân tử rất nhỏ so với kích thước của các phân tử.}
{các phân tử khí khác nhau va chạm vào thành bình tác dụng vào thành bình những lực bằng nhau.}
\loigiai{Trong điều kiện chuẩn về nhiệt độ và áp suất (điều kiện tiêu chuẩn áp suất 1 atm, nhiệt độ 273 K, thể tích 22,4 lít) thì số phân tử trong một đơn vị thể tích của các chất khí khác nhau là như nhau.}
\end{ex}

\dangbt{Phân biệt đặc điểm cấu trúc phân tử của chất ở thể rắn, lỏng, khí}
\begin{ex}
Áp suất của khí lên thành bình là do lực tác dụng
\choice
{lên một đơn vị diện tích thành bình}
{\True vuông góc lên một đơn vị diện tích thành bình.}
{lực tác dụng lên thành bình.}
{vuông góc lên toàn bộ diện tích thành bình.}
\loigiai{Áp suất là lực tác dụng vuông góc lên một đơn vị diện tích.}
\end{ex}

\dangbt{Phân biệt đặc điểm cấu trúc phân tử của chất ở thể rắn, lỏng, khí}
\begin{ex}
Phát biểu nào sau đây là sai khi nói về chất lỏng?
\choice
{\True Chất lỏng không có thể tích riêng xác định.}
{Các nguyên tử, phân tử cũng dao động quanh các vị trí cân bằng, nhưng những vị trí cân bằng này không cố định mà di chuyển.}
{Lực tương tác giữa các phân tử chất lỏng lớn hơn lực tương tác giữa các nguyên tử, phân tử chất khí và nhỏ hơn lực tương tác giữa các nguyên tử, phân tử chất rắn.}
{Chất lỏng không có hình dạng riêng mà có hình dạng của phần bình chứa nó.}
\loigiai{Chất lỏng có thể tích riêng xác định, nhưng không có hình dạng riêng.}
\end{ex}

\dangbt{Phân biệt đặc điểm cấu trúc phân tử của chất ở thể rắn, lỏng, khí}
\begin{ex}
Khi nói về các tính chất của chất khí, phát biểu đúng là
\choice
{bành trướng là chiếm một phần thể tích của bình chứa.}
{khi áp suất tác dụng lên một lượng khí tăng thì thể tích của khí tăng đáng kể.}
{\True chất khí có tính dễ nén.}
{chất khí có khối lượng riêng lớn so với chất rắn và chất lỏng.}
\loigiai{Chất khí có tính dễ nén vì khoảng cách giữa các phân tử khí lớn so với chất lỏng hay rắn.}
\end{ex}

\dangbt{Phân biệt đặc điểm cấu trúc phân tử của chất ở thể rắn, lỏng, khí}
\begin{ex}
Chuyển động nào sau đây là chuyển động của riêng các phân tử ở thể lỏng?
\choice
{Chuyển động hỗn loạn không ngừng.}
{Dao động xung quanh các vị trí cân bằng cố định.}
{Chuyển động hoàn toàn tự do.}
{\True Dao động xung quanh các vị trí cân bằng không cố định.}
\loigiai{Các phân tử ở thể lỏng dao động xung quanh các vị trí cân bằng không cố định.}
\end{ex}

\dangbt{Giải thích hiện tượng và ứng dụng liên quan đến sự hóa hơi}
\begin{ex}
Hai chất khí có thể trộn lẫn vào nhau tạo nên một hỗn hợp khí đồng đều là vì
(1). các phân tử khí chuyển động nhiệt.
(2). hai chất khí đã cho không có phản ứng hoá học với nhau.
(3). giữa các phân tử khí có khoảng trống.
\choice
{(1) và (2).}
{(2) và (3).}
{(3) và (1).}
{\True cả (1), (2) và (3).}
\loigiai{Hai chất khí có thể trộn lẫn vào nhau tạo nên một hỗn hợp khí đồng đều là do các phân tử khí chuyển động nhiệt, giữa các phân tử khí có khoảng trống và hai chất khí không phản ứng hóa học với nhau.}
\end{ex}

\dangbt{Phân biệt đặc điểm cấu trúc phân tử của chất ở thể rắn, lỏng, khí}
\begin{ex}
Gọi $n_r, n_l, n_k$ lần lượt là mật độ phân tử của một chất ở thể rắn, thể lỏng và thể khí. Thứ tự đúng là
\choice
{\True $n_r > n_l > n_k$}
{$n_r < n_l < n_k$}
{$n_r = n_l = n_k$}
{$n_r > n_k > n_l$}
\loigiai{Mật độ phân tử chất ở thể rắn lớn hơn ở thể lỏng và lớn hơn ở thể khí. Trừ nước, mật độ nước khi ở thể rắn (nước đá) nhỏ hơn mật độ nước khi ở thể lỏng.}
\end{ex}

\dangbt{Phân biệt đặc điểm cấu trúc phân tử của chất ở thể rắn, lỏng, khí}
\begin{ex}
Trong các yếu tố sau
I. Lực liên kết giữa các phân tử.
II. Khoảng cách giữa các phân tử.
III. Nhiệt độ của các phân tử.
IV. Mật độ của các phân tử.
Để phân biệt các trạng thái rắn, lỏng, khí ta không dựa vào yếu tố thứ
\choice
{II.}
{IV.}
{I.}
{\True III.}
\loigiai{Yếu tố nhiệt độ không quyết định đến trạng thái của một chất. Một chất có thể tồn tại ở trạng thái rắn, lỏng hoặc khí ở cùng một nhiệt độ tùy thuộc vào áp suất.}
\end{ex}

\dangbt{Phân biệt đặc điểm cấu trúc phân tử của chất ở thể rắn, lỏng, khí}
\begin{ex}
Trong các chất rắn, lỏng, khí, chất khó nén là
\choice
{\True chất rắn, chất lỏng.}
{chất khí, chất rắn.}
{chỉ có chất rắn.}
{chất khí, chất lỏng.}
\loigiai{Chất rắn, chất lỏng mật độ phân tử cao nên khó nén, còn chất khí do mật độ thấp nên dễ nén.}
\end{ex}

\dangbt{Giải thích hiện tượng và ứng dụng liên quan đến sự hóa hơi}
\begin{ex}
Khi nhiệt độ trong một bình tăng cao, áp suất của khối khí trong bình cũng tăng lên vì
\choice
{phân tử va chạm với nhau nhiều hơn.}
{số lượng phân tử tăng.}
{\True phân tử khí chuyển động nhanh hơn.}
{khoảng cách giữa các phân tử tăng.}
\loigiai{Khi nhiệt độ trong bình kín tăng cao, các phân tử chất khí sẽ chuyển động nhanh hơn, nên va chạm với thành bình nhiều hơn, tạo ra một lực lớn tác động vào thành bình làm áp suất trong bình tăng lên.}
\end{ex}

\dangbt{Phân biệt đặc điểm cấu trúc phân tử của chất ở thể rắn, lỏng, khí}
\begin{ex}
Tính chất nào sau đây không phải là tính chất của chất khí?
\choice
{\True Các phân tử chuyển động hỗn loạn xung quanh các vị trí cân bằng cố định.}
{Chất khí có tính bành trướng, luôn chiếm toàn bộ thể tích bình chứa.}
{Chất khí dễ nén hơn chất lỏng và chất rắn.}
{Các phân tử chuyển động hỗn loạn và không ngừng.}
\loigiai{Tính chất các phân tử chuyển động hỗn loạn xung quanh các vị trí cân bằng cố định là của chất rắn.}
\end{ex}

\dangbt{Phân biệt đặc điểm cấu trúc phân tử của chất ở thể rắn, lỏng, khí}
\begin{ex}
Các vật rắn giữ được hình dạng và thể tích của chúng là do loại lực nào sau đây?
\choice
{Lực hấp dẫn.}
{Lực ma sát.}
{\True Lực tương tác phân tử.}
{Lực hạt nhân.}
\loigiai{Các vật rắn giữ được hình dạng và thể tích của chúng là do lực tương tác phân tử.}
\end{ex}

\dangbt{Mô hình động học phân tử về cấu tạo chất và chuyển động Brown}
\begin{ex}
Điều nào sau đây là sai khi nói về cấu tạo chất?
\choice
{\True Các nguyên tử hay phân tử chuyển động càng nhanh thì nhiệt độ của vật càng thấp.}
{Các nguyên tử, phân tử chuyến động hỗn loạn không ngừng.}
{Các nguyên tử, phân tử đồng thời hút nhau và đẩy nhau.}
{Các chất được cấu tạo từ các nguyên tử, phân tử.}
\loigiai{Theo thuyết động học phân tử của chất khí thì chất khí được cấu tạo từ các phân tử riêng rẽ, các phân tử khí chuyển động hỗn loạn không ngừng, chuyển động này càng nhanh thì nhiệt độ của vật càng cao.}
\end{ex}

\dangbt{Phân biệt đặc điểm cấu trúc phân tử của chất ở thể rắn, lỏng, khí}
\begin{ex}
Phát biểu nào sau đây là đúng khi nói về các trạng thái rắn, lỏng, khí của vật chất?
\choice
{\True Chất khí không có hình dạng và thế tích xác định.}
{Chất lỏng không có thể tích riêng xác định.}
{Lực tương tác giữa các nguyên tử, phân tử trong chất rắn là rất yếu.}
{Trong chất lỏng các nguyên tử, phân tử dao động quanh vị trí cân bằng cố định.}
\loigiai{Lực liên kết giữa các phân tử chất khí nhỏ không đáng kể nên chúng chuyển động tự do về mọi phía. Do đó chất khí không có hình dạng và thể tích xác định.
Lực liên kết giữa các phân tử chất rắn giữ cho các nguyên tử, phân tử dao động quanh vị trí cân bằng cố định.
Trong chất lỏng các nguyên tử, phân tử dao động xung quanh các vị trí cân bằng nhưng vị trí này không cố định. Do đó chất lỏng có thể tích riêng xác định, nhưng có hình dạng phụ thuộc vào hình dạng của bình chứa.}
\end{ex}

\dangbt{Phân biệt đặc điểm cấu trúc phân tử của chất ở thể rắn, lỏng, khí}
\begin{ex}
Khẳng định nào sau đây là sai khi nói về cấu tạo chất?
\choice
{Các chất được cấu tạo từ các hạt riêng gọi là nguyên tử, phân tử.}
{\True Các nguyên tử, phân tử đứng sát nhau và giữa chúng không có khoảng cách.}
{Lực tương tác giữa các phân tử ở thể rắn lớn hơn lực tương tác giữa các phân tử ở thể lỏng và thể khí.}
{Các nguyên tử, phân tử chất lỏng dao động xung quanh các vị trí cân bằng không cố định.}
\loigiai{Các chất được cấu tạo từ các hạt riêng biệt gọi là nguyên tử, phân tử, giữa chúng có khoảng cách. Khoảng cách giữa các phân tử khác nhau ở các thể rắn, lỏng và khí. Lực tương tác giữa các phân tử xếp theo thứ tự tăng dần lần lượt ở thể khí, lỏng và rắn. Lực tương tác trong chất lỏng chưa đủ lớn để giữ các nguyên tử, phân tử không chuyển động phân tán ra xa nhau nên chúng dao động xung quanh các vị trí cân bằng không cố định.}
\end{ex}

\dangbt{Mô hình động học phân tử về cấu tạo chất và chuyển động Brown}
\begin{ex}
Xét các tính chất sau đây của các phân tử
(I) Chuyển động không ngừng.
(II) Tương tác với nhau bằng lực hút và lực đẩy.
(III) Khi chuyển động va chạm với nhau.
Các phân tử chất rắn, chất lỏng có cùng tính chất nào?
\choice
{\True (I) và (II).}
{(II) và (III).}
{(III) và (I).}
{(I), (II) và (III).}
\loigiai{Các phân tử của chất rắn, chất lỏng đều có cùng tính chất chuyển động không ngừng (dao động xung quanh các vị trí cân bằng) và tương tác với nhau bằng lực tương tác phân tử (lực của chất rắn lớn hơn lực của chất lỏng). Vị trí cân bằng của các phân tử, nguyên tử chất rắn hoàn toàn xác định và các nguyên tử phân tử này dao động với biên độ nhỏ nên không va chạm với nhau, còn chất lỏng có vị trí cân bằng luôn thay đổi nên có thể va chạm nhau khi chuyển động.}
\end{ex}

\dangbt{Mô hình động học phân tử về cấu tạo chất và chuyển động Brown}
\begin{ex}
Theo thuyết động học phân tử, các phân tử vật chất luôn chuyển động không ngừng. Thuyết này áp dụng cho
\choice
{chất khí.}
{\True chất rắn, lỏng và khí.}
{chất lỏng.}
{chất rắn.}
\loigiai{Thuyết động học phân tử áp dụng cho các chất rắn, lỏng, khí.}
\end{ex}

\dangbt{Phân biệt đặc điểm cấu trúc phân tử của chất ở thể rắn, lỏng, khí}
\begin{ex}
Các tính chất nào sau đây là tính chất của các phân tử chất lỏng?
\choice
{Không có thể tích xác định.}
{\True Hình dạng phụ thuộc bình chứa.}
{Lực tương tác phân tử yếu.}
{Khoảng cách giữa các phân tử rất lớn.}
\loigiai{Vị trí cân bằng của các phân tử chất lỏng có thể thay đổi được nên chất lỏng có hình dạng phụ thuộc vào bình chứa.}
\end{ex}

\dangbt{Phân biệt đặc điểm cấu trúc phân tử của chất ở thể rắn, lỏng, khí}
\begin{ex}
Trong các tính chất sau, tính chất nào không phải của chất khí?
\choice
{\True Có hình dạng cố định.}
{Chiếm toàn bộ thể tích của bình chứa.}
{Tác dụng lực lên mọi phần diện tích bình chứa.}
{Thể tích giảm đáng kể khi tăng áp suất.}
\loigiai{Lực tương tác giữa các phân tử khí rất yếu nên chất khí chuyển động hỗn loạn và không có hình dạng xác định.}
\end{ex}

\dangbt{Phân biệt đặc điểm cấu trúc phân tử của chất ở thể rắn, lỏng, khí}
\begin{ex}
Trong điều kiện chuẩn về nhiệt độ và áp suất thì
\choice
{\True số phân tử trong một đơn vị thể tích của các chất khí khác nhau là như nhau.}
{các phân tử của các chất khí khác nhau chuyển động với vận tốc như nhau.}
{khoảng cách giữa các phân tử rất nhỏ so với kích thước của các phân tử.}
{các phân tử khí khác nhau va chạm vào thành bình tác dụng vào thành bình những lực bằng nhau.}
\loigiai{Trong điều kiện chuẩn về nhiệt độ và áp suất (điều kiện tiêu chuẩn áp suất 1 atm, nhiệt độ 273 K, thể tích 22,4 lít) thì số phân tử trong một đơn vị thể tích của các chất khí khác nhau là như nhau.}
\end{ex}

\dangbt{Mô hình động học phân tử về cấu tạo chất và chuyển động Brown}
\begin{ex}
Tính chất nào sau đây không phải là của phân tử?
\choice
{Chuyển động không ngừng.}
{Giữa các phân tử có khoảng cách.}
{\True Có lúc đứng yên, có lúc chuyển động.}
{Chuyển động càng nhanh thì nhiệt độ của vật càng cao.}
\loigiai{Các phân tử luôn chuyển động không ngừng, không có lúc đứng yên.}
\end{ex}

\dangbt{Mô hình động học phân tử về cấu tạo chất và chuyển động Brown}
\begin{ex}
Thông tin nào sau đây là không đúng khi nói về khối lượng mol và thể tích mol của một chất?
\choice
{Khối lượng mol của một chất được đo bằng khối lượng của một mol chất ấy.}
{Thể tích mol của một chất được đo bằng thể tích của một mol chất ấy.}
{Ở điều kiện tiêu chuẩn (0 độ C và 1atm) thể tích mol của mọi chất khí đều bằng 22,4l.}
{\True Ở điều kiện tiêu chuẩn (0 độ C và 1atm) thể tích mol của các chất khí khác nhau thì khác nhau.}
\loigiai{Ở điều kiện tiêu chuẩn (0 độ C và 1atm) thể tích mol của mọi chất khí đều bằng 22,4l.}
\end{ex}

\dangbt{Phân biệt đặc điểm cấu trúc phân tử của chất ở thể rắn, lỏng, khí}
\begin{ex}
Điều nào sau đây là sai khi nói về chất lỏng?
\choice
{\True Chất lỏng không có thể tích riêng xác định}
{Các nguyên tử, phân tử cũng dao động quanh các vị trí cân bằng, nhưng những vị trí cân bằng này không cố định mà di chuyển}
{Lực tương tác giữa các phân tử chất lỏng lớn hơn lực tương tác giữa các nguyên tử, phân tử chất khí và nhỏ hơn lực tương tác giữa các nguyên tử, phân tử chất rắn}
{Chất lỏng không có hình dạng riêng mà có hình dạng của phần bình chứa nó.}
\loigiai{Chất lỏng có thể tích riêng xác định, nhưng không có hình dạng riêng.}
\end{ex}

\dangbt{Phân biệt đặc điểm cấu trúc phân tử của chất ở thể rắn, lỏng, khí}
\begin{ex}
Phát biểu nào sau đây là đúng khi nói về vị trí của các nguyên tử, phân tử trong chất rắn?
\choice
{\True Các nguyên tử, phân tử nằm ở những vị trí xác định và chỉ có thể dao động xung quanh các vị trí cân bằng này.}
{Các nguyên tử, phân tử nằm ở những vị trí cố định.}
{Các nguyên tử, phân tử không có vị trí cố định mà luôn thay đổi.}
{Các nguyên tử, phân tử nằm ở những vị trí cố định, sau một thời gian nào đó chúng lại chuyển sang một vị trí cố định khác.}
\loigiai{Trong chất rắn, các nguyên tử, phân tử nằm ở những vị trí xác định và chỉ có thể dao động xung quanh các vị trí cân bằng này.}
\end{ex}

\dangbt{Phân biệt đặc điểm cấu trúc phân tử của chất ở thể rắn, lỏng, khí}
\begin{ex}
Mô hình động học phân tử có thể dùng để lời giải thích cấu trúc của các chất rắn, chất lỏng và chất khí. Các chất được cấu tạo từ các nguyên tử, phân tử và chúng chuyển động không ngừng. Sự sắp xếp và tương tác khác nhau giữa các nguyên tử, phân tử sẽ tạo nên các tính chất khác nhau cho chất rắn, chất lỏng, chất khí.
\choiceTF
{Các chất rắn có hình dạng và thể tích xác định vì các phân tử cấu tạo nên nó đứng yên tại chỗ.}
{\True Các chất khí dễ nén hơn các chất lỏng vì khoảng cách trung bình giữa các phân tử chất khí lớn hơn khoảng cách trung bình giữa các phân tử chất lỏng.}
{\True Chất lỏng không có hình dạng xác định vì các phân tử dao động quanh những vị trí cân bằng có thể di chuyển được.}
{Tất cả các chất rắn đều có cấu trúc mạng tinh thể nên có hình dạng hình học xác định.} 
\loigiai{a. Phát biểu này sai. Các phân tử vẫn dao động quanh vị trí cân bằng cố định.
b. Phát biểu này đúng.
c. Phát biểu này đúng.
d. Phát biểu này sai. Chỉ có chất rắn kết tinh mới có cấu trúc mạng tinh thể, chất rắn vô định hình (thủy tinh, nhựa, sôcôla...) không có cấu trúc mạng tinh thể.}
\end{ex}

\dangbt{Phân biệt đặc điểm cấu trúc phân tử của chất ở thể rắn, lỏng, khí}
\begin{ex}
Theo mô hình động học phân tử thì các phân tử cấu tạo nên các vật có sự tương tác với nhau. Sự khác biệt về độ lớn của lực tương tác giữa các phân tử trong chất rắn, chất lỏng, chất khí dẫn đến sự khác nhau về cấu trúc của chúng.
\choiceTF
{\True Lực tương tác giữa các phân tử chất rắn mạnh hơn lực tương tác giữa các phân tử chất khí.}
{\True Lực tương tác giữa các phân tử chất rắn rất mạnh nên giữa các phân tử chỉ dao động quanh các vị trí cân bằng cố định, vì vậy chất rắn có hình dạng xác định.}
{Các phân tử chất lỏng tương tác với nhau rất yếu nên các phân tử chất lỏng có thể chuyển động tự do ra xa nhau, vì vậy chất lỏng không có hình dạng xác định.}
{\True Các chất khí có khoảng cách giữa các phân tử rất lớn so với kích thước của chúng nên lực tương tác giữa các phân tử gần như không đáng kể.}
\loigiai{a. Phát biểu này đúng. 
b. Phát biểu này đúng. 
c. Phát biểu này sai. Lực tương tác giữa các phân tử chất lỏng vẫn đủ để giữ các phân tử chất lỏng không chuyển động phân tán ra xa nhau.
d. Phát biểu này đúng.}
\end{ex}

\dangbt{Phân biệt đặc điểm cấu trúc phân tử của chất ở thể rắn, lỏng, khí}
\begin{ex}
Hình dưới đây mô tả chuyển động phân tử ở các thể khác nhau. Hình cầu là phân tử, mũi tên là hướng chuyển động của phân tử. 
\choiceTF
{\True Hình a mô tả chuyển động ở thể khí.}
{Hình b mô tả chuyển động ở thể lỏng.}
{Hình b mô tả chuyển động ở thể rắn.}
{\True Khoảng cách giữa các phân tử trong chất lỏng lớn hơn khoảng cách giữa các phân tử trong chất rắn và nhỏ hơn khoảng cách giữa các phân tử trong chất khí.} 
\loigiai{a. Phát biểu này đúng. 
b. Phát biểu này sai. Hình b mô tả chuyển động ở thể rắn. 	
c. Phát biểu này sai. Hình c mô tả chuyển động ở thể lỏng.
d. Phát biểu này đúng.}
\end{ex}
